\documentclass{article}

\usepackage{amsmath,amssymb}
\usepackage{xurl}
\usepackage[hidelinks]{hyperref}
% Shared commands for the notes in docs/paper-gaps/.
% Every note loads this file via % Shared commands for the notes in docs/paper-gaps/.
% Every note loads this file via % Shared commands for the notes in docs/paper-gaps/.
% Every note loads this file via \input{command}. Project-wide notation
% belongs here, never in per-note redefinitions.

% --- Project configuration ------------------------------------------------
\newcommand{\projectIssueBase}{https://github.com/LionSR/QICLean/issues/}
\newcommand{\projectPrBase}{https://github.com/LionSR/QICLean/pull/}
\newcommand{\projectDocBase}{https://sirui-lu.com/QICLean/docs/}
% ---------------------------------------------------------------------------

\DeclareUnicodeCharacter{2080}{\ifmmode{}_{0}\else\textsubscript{0}\fi}
\DeclareUnicodeCharacter{2081}{\ifmmode{}_{1}\else\textsubscript{1}\fi}
\DeclareUnicodeCharacter{2082}{\ifmmode{}_{2}\else\textsubscript{2}\fi}
\DeclareUnicodeCharacter{2099}{\ifmmode{}_{n}\else\textsubscript{n}\fi}

\newcommand{\Lean}{\textsc{Lean}}

% A formal declaration name, upright sans, breakable at underscores.
\ExplSyntaxOn
\NewDocumentCommand{\leanid}{m}{
  \ifmmode
    \textsf{\detokenize{#1}}
  \else
    \group_begin:
    \urlstyle{sf}
    \tl_set:Nn \l_tmpa_tl {#1}
    \tl_replace_all:Nnn \l_tmpa_tl {\_} {_}
    \exp_args:NV \path \l_tmpa_tl
    \group_end:
  \fi
}
\ExplSyntaxOff

% Traceability links; use in footnotes, not in the running prose.
\newcommand{\ghissue}[1]{\href{\projectIssueBase#1}{GitHub issue \##1}}
\newcommand{\ghpr}[1]{\href{\projectPrBase#1}{PR \##1}}
\newcommand{\doclink}[1]{\href{\projectDocBase#1.html}{\path{#1}}}

% Domain notation, using \providecommand so a note-local refinement stays
% possible.
\providecommand{\tr}{\operatorname{tr}}

% Dirac notation and the identity operator, as in blueprint/src/macros.
\usepackage{dsfont}
\providecommand{\ket}[1]{|#1\rangle}
\providecommand{\bra}[1]{\langle #1|}
\providecommand{\braket}[2]{\langle #1 | #2 \rangle}
\providecommand{\ketbra}[2]{|#1\rangle\!\langle #2|}
\providecommand{\Id}{\mathds{1}}
\providecommand{\one}{\mathds{1}}
\providecommand{\C}{\mathbb{C}}
\providecommand{\MN}[1]{M_{#1}(\C)}
\providecommand{\MD}{M_D(\C)}

% Verdict marker: \gapnote{<kind>}{<status>}. Kind is the policy
% classification (clarification, local-correction, scope-restriction,
% unfaithful, false-source, open-gap); status is open, wip (elimination
% underway), resolved, or historical. Machine-read by the site index;
% prints a verdict line.
\newcommand{\gapnote}[2]{\par\noindent\textbf{Verdict:} #1 (#2).\par}
. Project-wide notation
% belongs here, never in per-note redefinitions.

% --- Project configuration ------------------------------------------------
\newcommand{\projectIssueBase}{https://github.com/LionSR/QICLean/issues/}
\newcommand{\projectPrBase}{https://github.com/LionSR/QICLean/pull/}
\newcommand{\projectDocBase}{https://sirui-lu.com/QICLean/docs/}
% ---------------------------------------------------------------------------

\DeclareUnicodeCharacter{2080}{\ifmmode{}_{0}\else\textsubscript{0}\fi}
\DeclareUnicodeCharacter{2081}{\ifmmode{}_{1}\else\textsubscript{1}\fi}
\DeclareUnicodeCharacter{2082}{\ifmmode{}_{2}\else\textsubscript{2}\fi}
\DeclareUnicodeCharacter{2099}{\ifmmode{}_{n}\else\textsubscript{n}\fi}

\newcommand{\Lean}{\textsc{Lean}}

% A formal declaration name, upright sans, breakable at underscores.
\ExplSyntaxOn
\NewDocumentCommand{\leanid}{m}{
  \ifmmode
    \textsf{\detokenize{#1}}
  \else
    \group_begin:
    \urlstyle{sf}
    \tl_set:Nn \l_tmpa_tl {#1}
    \tl_replace_all:Nnn \l_tmpa_tl {\_} {_}
    \exp_args:NV \path \l_tmpa_tl
    \group_end:
  \fi
}
\ExplSyntaxOff

% Traceability links; use in footnotes, not in the running prose.
\newcommand{\ghissue}[1]{\href{\projectIssueBase#1}{GitHub issue \##1}}
\newcommand{\ghpr}[1]{\href{\projectPrBase#1}{PR \##1}}
\newcommand{\doclink}[1]{\href{\projectDocBase#1.html}{\path{#1}}}

% Domain notation, using \providecommand so a note-local refinement stays
% possible.
\providecommand{\tr}{\operatorname{tr}}

% Dirac notation and the identity operator, as in blueprint/src/macros.
\usepackage{dsfont}
\providecommand{\ket}[1]{|#1\rangle}
\providecommand{\bra}[1]{\langle #1|}
\providecommand{\braket}[2]{\langle #1 | #2 \rangle}
\providecommand{\ketbra}[2]{|#1\rangle\!\langle #2|}
\providecommand{\Id}{\mathds{1}}
\providecommand{\one}{\mathds{1}}
\providecommand{\C}{\mathbb{C}}
\providecommand{\MN}[1]{M_{#1}(\C)}
\providecommand{\MD}{M_D(\C)}

% Verdict marker: \gapnote{<kind>}{<status>}. Kind is the policy
% classification (clarification, local-correction, scope-restriction,
% unfaithful, false-source, open-gap); status is open, wip (elimination
% underway), resolved, or historical. Machine-read by the site index;
% prints a verdict line.
\newcommand{\gapnote}[2]{\par\noindent\textbf{Verdict:} #1 (#2).\par}
. Project-wide notation
% belongs here, never in per-note redefinitions.

% --- Project configuration ------------------------------------------------
\newcommand{\projectIssueBase}{https://github.com/LionSR/QICLean/issues/}
\newcommand{\projectPrBase}{https://github.com/LionSR/QICLean/pull/}
\newcommand{\projectDocBase}{https://sirui-lu.com/QICLean/docs/}
% ---------------------------------------------------------------------------

\DeclareUnicodeCharacter{2080}{\ifmmode{}_{0}\else\textsubscript{0}\fi}
\DeclareUnicodeCharacter{2081}{\ifmmode{}_{1}\else\textsubscript{1}\fi}
\DeclareUnicodeCharacter{2082}{\ifmmode{}_{2}\else\textsubscript{2}\fi}
\DeclareUnicodeCharacter{2099}{\ifmmode{}_{n}\else\textsubscript{n}\fi}

\newcommand{\Lean}{\textsc{Lean}}

% A formal declaration name, upright sans, breakable at underscores.
\ExplSyntaxOn
\NewDocumentCommand{\leanid}{m}{
  \ifmmode
    \textsf{\detokenize{#1}}
  \else
    \group_begin:
    \urlstyle{sf}
    \tl_set:Nn \l_tmpa_tl {#1}
    \tl_replace_all:Nnn \l_tmpa_tl {\_} {_}
    \exp_args:NV \path \l_tmpa_tl
    \group_end:
  \fi
}
\ExplSyntaxOff

% Traceability links; use in footnotes, not in the running prose.
\newcommand{\ghissue}[1]{\href{\projectIssueBase#1}{GitHub issue \##1}}
\newcommand{\ghpr}[1]{\href{\projectPrBase#1}{PR \##1}}
\newcommand{\doclink}[1]{\href{\projectDocBase#1.html}{\path{#1}}}

% Domain notation, using \providecommand so a note-local refinement stays
% possible.
\providecommand{\tr}{\operatorname{tr}}

% Dirac notation and the identity operator, as in blueprint/src/macros.
\usepackage{dsfont}
\providecommand{\ket}[1]{|#1\rangle}
\providecommand{\bra}[1]{\langle #1|}
\providecommand{\braket}[2]{\langle #1 | #2 \rangle}
\providecommand{\ketbra}[2]{|#1\rangle\!\langle #2|}
\providecommand{\Id}{\mathds{1}}
\providecommand{\one}{\mathds{1}}
\providecommand{\C}{\mathbb{C}}
\providecommand{\MN}[1]{M_{#1}(\C)}
\providecommand{\MD}{M_D(\C)}

% Verdict marker: \gapnote{<kind>}{<status>}. Kind is the policy
% classification (clarification, local-correction, scope-restriction,
% unfaithful, false-source, open-gap); status is open, wip (elimination
% underway), resolved, or historical. Machine-read by the site index;
% prints a verdict line.
\newcommand{\gapnote}[2]{\par\noindent\textbf{Verdict:} #1 (#2).\par}


\title{The Energy after Trimming in the Polynomial-Time Algorithm\\
    of Landau, Vazirani and Vidick}
\author{}
\date{2026-10-25}

\begin{document}
\maketitle
\gapnote{local-correction}{resolved}

\section{The source statement}

The Hamiltonian is $H=\sum_{i=1}^{n-1}H_i$ on a chain of $n$ qudits, where
$H_i$ acts on the qudits $i,i+1$ and satisfies $0\le H_i\le I$; the ground
energy is $\varepsilon_0$, the ground state is $\Gamma$, and the spectral gap
is $\varepsilon$.\footnote{The source is \cite{Landau2013PolynomialTime},
    version~1, source file \path{algo53.tex}: the Hamiltonian at line~283,
    the constant $c_\varepsilon$ at lines~288--291, Lemma
    \path{claim:energy-overlap} at lines~297--305, Lemma
    \path{claim:overlap} at lines~307--317, Lemma \path{l:2} with its proof
    at lines~347--359, and Lemma \path{l:leading} with its proof at
    lines~493--508.}
For a cut $(i,i+1)$, $\mathrm{trim}_B$ keeps the $B$ leading Schmidt terms
across the cut, and $B_\delta$ is the bond dimension of
Lemma~\path{constantbondapprox}: the normalized trimmed ground state
$u=\mathrm{trim}_{B_\delta}\Gamma/\|\mathrm{trim}_{B_\delta}\Gamma\|$
satisfies $|\langle\Gamma,u\rangle|\ge1-\delta$. Lemma~\path{l:2} reads as
follows.

\begin{quote}
    Let $\delta>0$ be such that $\delta(1+1/\varepsilon)\le\frac12$ and $w$ a
    vector with energy no larger than $\varepsilon_0+\delta$. Then
    $v:=\mathrm{trim}_{B_\delta}w/\|\mathrm{trim}_{B_\delta}w\|$ has energy no
    larger than $\varepsilon_0+6\sqrt\delta$.
\end{quote}

The proof (lines~350--358) runs as follows. Lemma
\path{claim:energy-overlap} gives $|\langle\Gamma,w\rangle|\ge
1-\delta/\varepsilon$, and Lemma \path{claim:overlap} is quoted as giving
\begin{align}
    |\langle w,u\rangle| & \ge1-\delta(1+1/\varepsilon)=1-\delta'.
    \label{eq:lvv13_trimming_energy_printed_overlap}
\end{align}
The Eckart--Young theorem is quoted as giving
$\|\mathrm{trim}_{B_\delta}w\|^2\ge1-\delta'$. With
$w_0=\mathrm{trim}_{B_\delta}w$ and $w_1=w-w_0$, the proof writes
\begin{align}
    \langle w,Hw\rangle & =\langle w_0,(H-H_i)w_0\rangle
    +\langle w_1,(H-H_i)w_1\rangle+\langle w,H_iw\rangle,
    \label{eq:lvv13_trimming_energy_split}
\end{align}
bounds $|\langle w_0,H_iw_0\rangle-\langle w,H_iw\rangle|\le2\|w_1\|\le
2\sqrt{\delta'}$ and $\langle w_1,(H-H_i)w_1\rangle\ge(\varepsilon_0-1)
\|w_1\|^2\ge(\varepsilon_0-1)\delta'$, and concludes
\begin{align}
    \langle w_0,Hw_0\rangle & \le\varepsilon_0(1-\delta')+\delta+2\sqrt\delta,
    \label{eq:lvv13_trimming_energy_printed_w0}                        \\
    \langle v,Hv\rangle     & \le\varepsilon_0
    +\frac{\delta+2\sqrt\delta}{1-\delta'}
    \le\varepsilon_0+6\sqrt\delta\quad\text{if }\delta'\le\tfrac12.
    \label{eq:lvv13_trimming_energy_printed}
\end{align}

\section{Three slips}

\paragraph{The factor two of the overlap lemma.} Lemma \path{claim:overlap}
states that $|\langle v,w\rangle|\ge1-\delta_1$ and $|\langle v',w\rangle|\ge
1-\delta_2$ give $|\langle v,v'\rangle|\ge1-2(\delta_1+\delta_2)$. Applied
with $\delta_1=\delta/\varepsilon$ and $\delta_2=\delta$ it gives
\begin{align}
    |\langle w,u\rangle| & \ge1-2\delta(1+1/\varepsilon)=1-\eta,
     & \eta & =2\delta(1+1/\varepsilon),
    \label{eq:lvv13_trimming_energy_overlap}
\end{align}
not (\ref{eq:lvv13_trimming_energy_printed_overlap}).

\paragraph{The square in Eckart--Young.} Since $u$ is a unit vector of
Schmidt rank at most $B_\delta$, the Eckart--Young theorem gives
$\|\mathrm{trim}_{B_\delta}w\|\ge|\langle u,w\rangle|$, a bound on the norm
and not on its square. With $t=\|w_1\|^2$ and
$\|w_0\|^2+\|w_1\|^2=1$, the overlap (\ref{eq:lvv13_trimming_energy_overlap})
gives
\begin{align}
    1-t=\|w_0\|^2 & \ge(1-\eta)^2\ge1-2\eta,
     & \text{so }t & \le2\eta=4\delta(1+1/\varepsilon)
    \label{eq:lvv13_trimming_energy_mass}
\end{align}
when $\eta\le1$; for $\eta>1$ the bound $t\le1<2\eta$ is trivial.

\paragraph{The dropped term and the square root.} The lower bound
$\langle w_1,(H-H_i)w_1\rangle\ge(\varepsilon_0-1)t$ inserted into
(\ref{eq:lvv13_trimming_energy_split}) gives
$\langle w_0,Hw_0\rangle\le\varepsilon_0(1-t)+\delta+t+2\sqrt t$. The
display (\ref{eq:lvv13_trimming_energy_printed_w0}) drops the term $t$,
replaces $(\varepsilon_0-1)t$ by $(\varepsilon_0-1)\delta'$ although only
$t\le\delta'$ is available, which gives a lower bound only when
$\varepsilon_0\le1$, and writes $2\sqrt\delta$ for the
$2\sqrt{\delta'}$ of the line before. Since $\delta'\ge\delta$, the last
change is not an upper bound either.

\section{The corrected bound}

Let $X=H-H_i$, the sum of the terms on the two sides of the cut. The
trimming projection $P$ commutes with the reduced density matrix of $w$ on
the left factor, so $X$ has no cross terms between $w_0=(P\otimes I)w$ and
$w_1=((I-P)\otimes I)w$:
\begin{align}
    \langle w,Xw\rangle & =\langle w_0,Xw_0\rangle+\langle w_1,Xw_1\rangle.
    \label{eq:lvv13_trimming_energy_no_cross}
\end{align}
This is the identity (\ref{eq:lvv13_trimming_energy_split}). Expanding
$\langle w,H_iw\rangle$ along $w=w_0+w_1$ and keeping
$\langle w_1,H_iw_1\rangle$ gives the exact identity
\begin{align}
    \langle w_0,Hw_0\rangle & =\langle w,Hw\rangle-\langle w_1,Hw_1\rangle
    -2\operatorname{Re}\langle w_0,H_iw_1\rangle.
    \label{eq:lvv13_trimming_energy_identity}
\end{align}
Now $\langle w_1,Hw_1\rangle\ge\varepsilon_0t$, and since $0\le H_i\le I$
the Cauchy--Schwarz inequality for the form of $H_i$ gives
$|\operatorname{Re}\langle w_0,H_iw_1\rangle|\le\|w_0\|\|w_1\|\le\sqrt t$.
Hence
\begin{align}
    \langle w_0,Hw_0\rangle & \le\varepsilon_0(1-t)+\delta+2\sqrt t,
     & \langle v,Hv\rangle & \le\varepsilon_0+\frac{\delta+2\sqrt t}{1-t}
    \label{eq:lvv13_trimming_energy_corrected}
\end{align}
for $t<1$, and with (\ref{eq:lvv13_trimming_energy_mass}), if $2\eta<1$,
\begin{align}
    \langle v,Hv\rangle & \le\varepsilon_0
    +\frac{\delta+2\sqrt{2\eta}}{1-2\eta}
    =\varepsilon_0+\frac{\delta+4\sqrt{\delta(1+1/\varepsilon)}}
    {1-4\delta(1+1/\varepsilon)}.
    \label{eq:lvv13_trimming_energy_final}
\end{align}
Under $\delta(1+1/\varepsilon)\le\frac18$ this is at most
$\varepsilon_0+2\delta+8\sqrt{\delta(1+1/\varepsilon)}$. The positivity
$H_i\ge0$ is what removes the term $t$ of the printed argument; with only
$\|H_i\|\le1$ the bound carries an extra $t/(1-t)$.

\section{The printed bound}

The bound (\ref{eq:lvv13_trimming_energy_final}) grows as
$\sqrt{\delta/\varepsilon}$ for small $\varepsilon$, while the printed bound
$\varepsilon_0+6\sqrt\delta$ is uniform in $\varepsilon$. The hypothesis
$\delta(1+1/\varepsilon)\le\frac12$ of the lemma gives only $\eta\le1$, which
does not even make $t\le\frac12$. The printed argument therefore does not
establish $\varepsilon_0+6\sqrt\delta$. No counterexample is known to us,
and this note does not claim that the printed bound fails; it records that
the formalized statement is (\ref{eq:lvv13_trimming_energy_final}).

\section{Effect downstream}

The constant $c_\varepsilon=(\varepsilon/169)^2$ (line~290) is chosen so
that $c_\varepsilon(1+1/\varepsilon)\le\frac12$,
$14\sqrt{c_\varepsilon}/\varepsilon<\frac1{12}$ and
$84c_\varepsilon/\varepsilon<\frac12$. Lemma \path{l:within} applies
Lemma \path{l:2} with $\delta=c_\varepsilon$ to obtain a vector of bond
dimension $B_{c_\varepsilon}$ with energy at most
$\varepsilon_0+6\sqrt{c_\varepsilon}$, and Lemma \path{l:leading}
(lines~493--508) uses this energy $\varepsilon_0+E$ in the form
$\operatorname{tr}(\sigma'H)\le\varepsilon_0+E+c_\varepsilon$, then a
Markov bound with threshold $2(E+c_\varepsilon)$, which must be smaller than
$\varepsilon/12$; the source has $E+c_\varepsilon\le7\sqrt{c_\varepsilon}$
and threshold $14\sqrt{c_\varepsilon}$.

With the corrected bound and $\delta=c_\varepsilon$,
$E\le2c_\varepsilon+8\sqrt{c_\varepsilon(1+1/\varepsilon)}$ once
$c_\varepsilon(1+1/\varepsilon)\le\frac18$, and the requirement becomes
\begin{align}
    2\Bigl(3c_\varepsilon+8\sqrt{c_\varepsilon(1+1/\varepsilon)}\Bigr)
     & <\frac{\varepsilon}{12}.
    \label{eq:lvv13_trimming_energy_requirement}
\end{align}
For $\varepsilon\le1$ one has $1+1/\varepsilon\le2/\varepsilon$, and the
choice $c_\varepsilon=\kappa\varepsilon^3$ turns
(\ref{eq:lvv13_trimming_energy_requirement}) into
$\varepsilon(6\kappa\varepsilon^2+16\sqrt{2\kappa})<\varepsilon/12$, which
holds for $\kappa=10^{-5}$ since $6\cdot10^{-5}+16\sqrt{2\cdot10^{-5}}
<0.0717<\frac1{12}$. For $\varepsilon>1$ the constant $10^{-5}$ works
directly, since then $1+1/\varepsilon<2$ and the left side of
(\ref{eq:lvv13_trimming_energy_requirement}) is below $0.072<1/12$. So
\begin{align}
    c_\varepsilon & =10^{-5}\min(\varepsilon,1)^3
    \label{eq:lvv13_trimming_energy_new_constant}
\end{align}
satisfies (\ref{eq:lvv13_trimming_energy_requirement}),
$c_\varepsilon(1+1/\varepsilon)\le\frac18$, $c_\varepsilon\le\varepsilon$ (needed
in Lemma \path{claim:energy-overlap}) and $84c_\varepsilon/\varepsilon<\frac12$.
The quadratic $c_\varepsilon=(\varepsilon/169)^2$ does not suffice for
small $\varepsilon$: the left side of
(\ref{eq:lvv13_trimming_energy_requirement}) is then of order
$\sqrt\varepsilon$, not $\varepsilon$. The change from $\varepsilon^2$ to
$\varepsilon^3$ enters the bond dimension $B_{c_\varepsilon}=
\exp(O(\varepsilon^{-1}\log^3d\,\log(1/c_\varepsilon)))$ only through
$\log(1/c_\varepsilon)$, which changes by a constant factor, so the
polynomial running time of the algorithm is unaffected.

\section{Formalization}

The formalization works with one cut, a bipartite space whose factors
carry the operators $A$ and $B$, and $H=A\otimes I+I\otimes B+K$ with
$H\ge\varepsilon_0$ and $0\le K\le I$; the operator $K$ plays the role of
$H_i$. The trimming projection is
\leanid{Matrix.IsTrimmingProjection} and the trimmed vector is
\leanid{Matrix.leftApply}. The orthogonality $\|w_0\|^2+\|w_1\|^2=\|w\|^2$
is \leanid{Matrix.IsTrimmingProjection.norm_leftApply_sq_add_norm_sub_sq}, and
(\ref{eq:lvv13_trimming_energy_no_cross}) is
\leanid{Matrix.IsTrimmingProjection.inner_toEuclideanLin_eq_add}. The first
bound of (\ref{eq:lvv13_trimming_energy_corrected}) is
\leanid{Matrix.IsTrimmingProjection.re_inner_leftApply_le}, the second is
\leanid{Matrix.IsTrimmingProjection.re_inner_normalized_le}, and its form for
$t\le\frac12$ is
\leanid{Matrix.IsTrimmingProjection.re_inner_normalized_le_of_le_half}. The
mass bound (\ref{eq:lvv13_trimming_energy_mass}) is
\leanid{Matrix.IsTrimmingProjection.norm_sub_leftApply_sq_le}, and
(\ref{eq:lvv13_trimming_energy_final}), stated through an overlap
$|\langle u,w\rangle|\ge1-\eta$ with a unit vector $u$ of Schmidt rank at
most the trimming radius, is
\leanid{Matrix.IsTrimmingProjection.re_inner_normalized_le_of_norm_inner}.
The overlap lemma is \leanid{one_sub_two_mul_add_le_norm_inner}. The printed
bound $\varepsilon_0+6\sqrt\delta$ is not formalized.

\bibliographystyle{alpha}
\bibliography{references}

\end{document}
